\section*{Step 260: Zero-Density Extension}

\paragraph{Carrier.}
For an $L$-function $L$ in the Selberg/automorphic class, define
\[
N_L(\sigma,T)=\#\{\rho=\beta+i\gamma:\ L(\rho)=0,\ \beta\geq \sigma,\ |\gamma|\leq T\}.
\]
The zero-density residual is
\[
\Xi_{\mathrm{density},L}(\sigma;T)
  = \frac{\log \max(1,N_L(\sigma,T))}{\log T}-f_0(\sigma).
\]
For the Riemann zeta function the canonical density-hypothesis exponent is
$f_0(\sigma)=2(1-\sigma)$.

\paragraph{Classification.}
Zero-density is a quantitative zero-count residual.  It is not zero-location
closure, not zero or coefficient correlation, and not a critical-line growth or
moment asymptotic.  Therefore it is typed as a separate Selberg-class extension,
parallel to SCDG, Cross-Correlation, and Subconvexity.

\paragraph{Candidate Statement.}
\emph{Selberg-Class Zero-Density Extension.} Beyond zero location,
zero/coefficient correlations, and critical-line growth, Selberg-class cascades
admit quantitative zero-count residuals in right-half rectangles.  Closure
conditions are upper bounds
\[
N_L(\sigma,T)\ll_{\epsilon,L}T^{f_0(\sigma)+\epsilon}.
\]
The zeta density hypothesis $f_0(\sigma)=2(1-\sigma)$ is the canonical target.

\paragraph{Evidence.}
The typed family is supported by three instances: zeta zero-density estimates
(Selberg, Bombieri, Huxley), Dirichlet-family average density
(Bombieri--Vinogradov / large sieve), and automorphic-family density estimates.

\paragraph{Nonclaim.}
This step does not prove the density hypothesis, Lindelof, GRH, or RH.  RH
implies density trivially for $\sigma>1/2$; Lindelof-type growth control yields
density consequences, but no converse is asserted.
